\documentclass[11pt]{article}

\usepackage[a4paper,margin=1in]{geometry}
\usepackage{amsmath,amssymb}
\usepackage{hyperref}
\usepackage{setspace}

\setstretch{1.12}
\hypersetup{colorlinks=true,citecolor=blue,linkcolor=blue,urlcolor=blue}
\urlstyle{same}

\title{\textbf{BABAPPAi: an applicability-aware framework around a frozen BABAPPA$\Omega$ backbone for branch--site identifiability diagnostics}}

\author{
Krishnendu Sinha\\
Assistant Professor, Department of Zoology, Jhargram Raj College, Jhargram 721507, India
}

\date{}

\begin{document}
\maketitle

\begin{abstract}
Branch--site analyses are often interpreted as evidence for episodic positive selection, but a prior question is whether episodic structure is \emph{recoverable} from the available alignment. BABAPPAi targets this identifiability question and separates it from adaptation claims while using a fixed neural inference backbone: the canonical frozen BABAPPA$\Omega$ model artifact. We lock the raw dispersion statistic as $D_{\mathrm{obs}}=\mathrm{Var}_{\mathrm{ddof}=1}(W_{0.05}(\overline{\ell}_i))$, where $\overline{\ell}_i$ is the site-level mean branch logit and $W_{0.05}$ denotes two-sided 5\% winsorization, and define raw $ \mathrm{EII}_z=(D_{\mathrm{obs}}-\mu_0)/\sigma_0 $ under matched-neutral Monte Carlo calibration. We then calibrate raw EII on held-out simulator truth \emph{conditionally on applicability}: $c\mathrm{EII}_{\mathrm{gene}}=P(I_{\mathrm{gene}}=1\mid \text{data},A_{\mathrm{gene}}=\text{applicable},\texttt{ceii\_v2})$ and $c\mathrm{EII}_{\mathrm{site}}=P(I_{\mathrm{site}}=1\mid \text{data},A_{\mathrm{site}}=\text{applicable},\texttt{ceii\_v2})$. Outside validated support, BABAPPAi abstains from calibrated probability output and reports explicit applicability diagnostics instead. Recoverability targets are truth-aware composites of site enrichment, site/branch rank agreement, burden alignment, and replicate-stability terms, with pre-registered binary labels for gene- and site-scale identifiability. Calibration uses a two-stage design (applicability gate + monotone score calibration) on a dedicated calibration split with separate test and regime-holdout (OOD) evaluation. Matched-neutral significance remains a distinct inferential layer via empirical $p_{\mathrm{emp}}$ and BH-adjusted $q_{\mathrm{emp}}$. In the current locked benchmark, gene-scale calibration was more stable than site-scale calibration, and OOD degradation was strongest for site localization under high nuisance regimes. The contribution is therefore a stable fixed-model inference stack plus honest post-inference calibration/abstention semantics that complement, but do not replace, classical branch--site hypothesis tests.
\end{abstract}

\noindent\textbf{Keywords:} episodic selection; branch--site evolution; likelihood-free inference; simulation-based inference; mutation--selection models; phylogenomics

\section*{Key points}
\begin{itemize}
    \item BABAPPA$\Omega$ is the canonical frozen inference model; BABAPPAi is the software framework around this fixed backbone.
    \item BABAPPAi targets the \emph{identifiability} of branch--site episodic signal rather than direct hypothesis testing for positive selection.
    \item The framework combines forward simulation under mutation--selection dynamics with likelihood-free neural inference trained on exact latent labels.
    \item Aggregate burden summaries are more recoverable than site-wise signals, indicating that episodic structure may become measurable at gene scale before it is precisely localizable.
    \item The software reports raw dispersion diagnostics ($\mathrm{EII}_{z,\mathrm{raw}}$, $\mathrm{EII}_{01,\mathrm{raw}}$) universally, and reports calibrated identifiability probabilities ($c\mathrm{EII}_{\mathrm{gene}}$, $c\mathrm{EII}_{\mathrm{site}}$) only when applicability criteria are satisfied.
    \item Matched-neutral significance (\texttt{p\_emp}, \texttt{q\_emp}, \texttt{significant\_bool}) is a separate inferential layer and is not interchangeable with calibrated identifiability probability.
\end{itemize}

\section{Introduction}

Detecting adaptive evolution in protein-coding genes is a central objective of comparative genomics. Codon-based likelihood methods, especially implementations in PAML and HyPhy, have dominated this area because they provide explicit statistical models for nonsynonymous and synonymous substitution processes \cite{Yang2007,Murrell2012,KosakovskyPond2020}. Within this framework, branch--site models are particularly attractive because they allow elevated nonsynonymous rates at a subset of codons along prespecified lineages, matching the biological expectation that adaptation is often sparse in both time and sequence space \cite{AnisimovaYang2007,YangDosReis2011,KosakovskyPond2011}.

Yet classical branch--site inference addresses a narrower target than is often assumed. In practice, these methods do not simply ask whether episodic selection exists in some biological sense; they ask whether the data support that conclusion under a specific codon-model parameterization. This distinction matters because limited alignment length, shallow or excessive divergence, recombination, and alignment error can all degrade statistical recoverability. Earlier simulation studies showed that the performance of likelihood-ratio procedures depends strongly on sequence divergence, model misspecification and alignment quality \cite{Anisimova2001,Wong2004,FletcherYang2010,JordanGoldman2012}. Recombination is especially problematic because most codon-model analyses assume a single genealogy across the alignment, whereas real coding sequences may be mosaics of distinct local histories \cite{Anisimova2003}. As a result, the absence of statistical support is not equivalent to the absence of biological signal.

This gap motivates a different inferential question: \emph{is episodic signal recoverable from the alignment at all, and at what scale?} Recent advances in simulation-based and likelihood-free inference make this question tractable. Neural networks have already shown substantial utility in population-genetic inference because they can learn informative mappings from complex sequence patterns without requiring closed-form likelihoods \cite{SheehanSong2016,Flagel2019,Cranmer2020}. In evolutionary genomics, this is especially useful when the simulator can encode richer biological structure than the downstream analytic model.

BABAPPAi was developed from this perspective. Instead of treating episodic selection as a shift in codon-model $\omega$ classes, the framework represents it as branch- and site-specific perturbation of latent codon-fitness landscapes within a mutation--selection process \cite{HalpernBruno1998}. A neural model is then trained on simulator-generated alignments with exact latent labels to infer branch--site deviation scores from observed sequence patterns. These scores are aggregated into burden summaries and an Episodic Identifiability Index (EII), which measures whether the alignment contains enough structured variation for episodic interpretation to be statistically meaningful. BABAPPAi therefore addresses a complementary problem to classical branch--site testing: not whether a codon-model null can be rejected, but whether episodic signal is recoverable from the data, and at what scale. In this manuscript the inference engine is fixed: the canonical frozen BABAPPA$\Omega$ model is used unchanged across all experiments, while BABAPPAi provides the operational software, significance layer, applicability checks, and calibration/reporting stack around that fixed model.

\section{Materials and methods}

\subsection{Framework overview}

BABAPPAi consists of three components: (i) a forward simulator of coding-sequence evolution under branch-specific mutation--selection dynamics, (ii) a fixed neural inference model (BABAPPA$\Omega$) trained on simulator-generated alignments with exact latent labels, and (iii) a calibration layer that summarizes branch--site scores into aggregate burden measures and the Episodic Identifiability Index. The framework is likelihood-free at inference time: it does not perform gene-wise likelihood optimization, branch-wise $\omega$ estimation, or likelihood-ratio testing.

\subsection{Forward simulation under mutation--selection dynamics}

Evolution is simulated on the 61-sense-codon state space under the standard genetic code, excluding stops. At each site $i \in \{1,\ldots,L\}$, a latent codon-fitness landscape is defined as
\[
F_i:\mathcal{C}\rightarrow \mathbb{R},
\]
where $F_i(c)$ denotes the fitness of codon $c$ at site $i$. Baseline landscapes are sampled as
\[
F_i^{(0)}(c)\sim \mathcal{N}(0,\sigma_F^2),
\]
thereby allowing a continuum of weakly deleterious, nearly neutral and mildly favourable states without imposing discrete codon-model classes.

To introduce gene-wide codon-usage heterogeneity, a codon-specific offset is added:
\[
F_i(c)=F_i^{(0)}(c)+\gamma \log\!\left(\frac{p_c}{\bar p}\right),
\]
where $p_c$ is sampled from a Dirichlet distribution over the codon state space, $\bar p$ is the mean codon frequency, and $\gamma$ modulates the contribution of codon-usage preference.

Episodic selection is generated as a sparse, branch- and site-specific perturbation of the latent fitness landscape. For selected branch--site pairs $(i,b)$,
\[
F_{i,b}(c)=F_i(c)+\Delta F_{i,b}(c),
\qquad
\Delta F_{i,b}(c)\sim \mathcal{N}(0,\sigma_{\mathrm{epi}}^2).
\]
Only a sparse subset of branch--site pairs receives perturbation, preserving the expectation that episodic adaptation is localized in both lineage and position. To approximate local genealogical heterogeneity, the alignment can be partitioned into contiguous blocks that act as coarse recombination-aware segments.

Single-nucleotide mutations follow an HKY-type scheme with transition--transversion bias \cite{Hasegawa1985}. If codons $c$ and $c'$ differ by exactly one nucleotide,
\[
\mu_{c\rightarrow c'}=
\begin{cases}
\kappa\mu, & \text{if the change is a transition},\\
\mu, & \text{if the change is a transversion},
\end{cases}
\]
where $\mu$ is the baseline mutation rate and $\kappa$ is the transition bias parameter. For branch $b$ and site $i$, the selection coefficient of a proposed change is
\[
s_{i,b}(c\rightarrow c')=F_{i,b}(c')-F_{i,b}(c).
\]
Substitutions are then realized under a Halpern--Bruno-style mutation--selection process \cite{HalpernBruno1998}, with continuous-time trajectories simulated using a Gillespie algorithm \cite{Gillespie1977}.

Because episodic perturbations are specified explicitly, exact latent labels are available. For each site--branch pair $(i,b)$, the simulator records an episodic indicator $Y_{i,b}\in\{0,1\}$. These indicators are aggregated into site- and branch-level burdens:
\[
B_i^{\mathrm{site}}=\frac{1}{|\mathcal{B}|}\sum_{b\in\mathcal{B}} Y_{i,b},
\qquad
B_b^{\mathrm{branch}}=\frac{1}{L}\sum_{i=1}^{L} Y_{i,b},
\]
where $\mathcal{B}$ is the set of branches in the tree. These burden quantities provide ground truth for training and benchmarking.

\subsection{Neural inference}

Each training example consists of a codon alignment together with branch-conditioned features derived from parent--child codon transitions, branch lengths and positional context. The model is trained in two stages. First, it is pretrained on a large corpus of simulated alignments spanning neutral, weakly structured and strongly episodic regimes. Second, it is fine-tuned against aggregate burden targets to align network outputs with site- and branch-level latent burden.

Let $\ell_{i,b}$ denote the site--branch logit predicted by the model, and let $\sigma(\cdot)$ be the logistic function. Site- and branch-level predicted burdens are
\[
\hat B_i^{\mathrm{site}}=\frac{1}{|\mathcal{B}|}\sum_{b\in\mathcal{B}}\sigma(\ell_{i,b}),
\qquad
\hat B_b^{\mathrm{branch}}=\frac{1}{L}\sum_{i=1}^{L}\sigma(\ell_{i,b}).
\]
Training minimizes a composite burden-aligned loss,
\[
\mathcal{L}_{\mathrm{total}}=\lambda_s\mathcal{L}_{\mathrm{site}}+\lambda_b\mathcal{L}_{\mathrm{branch}},
\]
where $\lambda_s$ and $\lambda_b$ are fixed loss weights. After training, model parameters are frozen and used as an immutable inference artifact. The release analysed here does not retrain this model; all updates are confined to post-inference interpretation layers.

At inference time, BABAPPAi returns branch--site deviation scores, site-level burdens, branch-level burdens, and gene-level summaries obtained by averaging site scores within each coding region. These outputs are diagnostic and relative; they are not posterior probabilities of positive selection and they do not replace formal codon-model likelihood statistics.

\subsection{Episodic Identifiability Index}

We lock one dispersion definition for all validation and release outputs:
\[
D_{\mathrm{obs}}=\mathrm{Var}_{\mathrm{ddof}=1}\!\left(W_{0.05}(\overline{\ell}_i)\right), \qquad
\overline{\ell}_i=\frac{1}{|\mathcal{B}|}\sum_{b\in\mathcal{B}}\ell_{i,b},
\]
where $\ell_{i,b}$ is the model logit for site $i$ and branch $b$, and $W_{0.05}$ is two-sided 5\% winsorization across sites. This removes prior ambiguity about equivalent dispersion summaries.

Matched-neutral calibration is also fixed: for each observed alignment, neutral replicates are matched on tree topology/size and alignment length, and the same $D$ statistic is recomputed to obtain $(\mu_0,\sigma_0)$ and neutral replicate values $D_0^{(1)},\ldots,D_0^{(M)}$.

Raw EII is defined as
\[
\mathrm{EII}_{z,\mathrm{raw}}=\frac{D_{\mathrm{obs}}-\mu_0}{\sigma_0}, \qquad
\mathrm{EII}_{01,\mathrm{raw}}=\sigma\!\left(\mathrm{EII}_{z,\mathrm{raw}}\right).
\]
Raw EII is an effect-size-style dispersion magnitude, not a probability.

To separate ``signal exists'' from ``signal recoverable'', we define truth-aware recoverability targets on held-out simulations:
\[
R_{\mathrm{gene}}=0.45\cdot \widetilde{\rho}_{\mathrm{branch}}
+0.35\cdot S_{\mathrm{burden}}
+0.20\cdot S_{\mathrm{branch,stab}},
\]
\[
R_{\mathrm{site}}=0.45\cdot E@k_{\mathrm{site}}
+0.35\cdot \widetilde{\rho}_{\mathrm{site}}
+0.20\cdot S_{\mathrm{site,stab}},
\]
where $\widetilde{\rho}=(\rho+1)/2$ is normalized Spearman rank agreement, $S_{\mathrm{burden}}$ is burden-alignment score, $E@k_{\mathrm{site}}$ is site enrichment at truth-derived $k$, and $S_{\mathrm{stab}}$ terms summarize replicate robustness.

Pre-registered binary identifiability labels are
\[
I_{\mathrm{gene}}=\mathbf{1}[R_{\mathrm{gene}}\ge \tau_{\mathrm{gene}}], \quad
\tau_{\mathrm{gene}}=0.42,
\]
\[
I_{\mathrm{site}}=\mathbf{1}[R_{\mathrm{site}}\ge \tau_{\mathrm{site}}], \quad
\tau_{\mathrm{site}}=0.45.
\]
By design, $I_{\mathrm{site}}$ is stricter in the benchmarked domain.

Calibration in \texttt{ceii\_v2} is two-stage. Stage 1 is an applicability function $A(x)$ over covariates $x$ (minimum: taxa count and gene length; optional additional run covariates):
\[
A(x)\in\{\texttt{in\_domain},\texttt{near\_boundary},\texttt{out\_of\_domain}\}.
\]
Only \texttt{in\_domain} samples are calibrated by default. Stage 2 fits a monotone mapping on a learned linear score from $(\mathrm{EII}_{z,\mathrm{raw}},X)$:
\[
c\mathrm{EII}_{\mathrm{gene}}=P(I_{\mathrm{gene}}=1\mid \mathrm{EII}_{z,\mathrm{raw}},X,A_{\mathrm{gene}}=\texttt{in\_domain},\texttt{ceii\_v2}),
\]
\[
c\mathrm{EII}_{\mathrm{site}}=P(I_{\mathrm{site}}=1\mid \mathrm{EII}_{z,\mathrm{raw}},X,A_{\mathrm{site}}=\texttt{in\_domain},\texttt{ceii\_v2}).
\]
When $A\neq \texttt{in\_domain}$ (or applicability is insufficient), BABAPPAi abstains:
\[
c\mathrm{EII}_{\mathrm{gene}}=\varnothing,\qquad c\mathrm{EII}_{\mathrm{site}}=\varnothing,
\]
and emits explicit applicability metadata (\texttt{applicability\_score}, \texttt{applicability\_status}, \texttt{calibration\_unavailable\_reason}, nearest supported regime, and distance-to-domain diagnostics). Bootstrap resampling provides calibration-curve uncertainty.

\subsection{Applicability envelope and calibration abstention}

The release implementation exposes applicability as a first-class selective-prediction layer. For each run, BABAPPAi reports \texttt{applicability\_score}, \texttt{within\_applicability\_envelope}, \texttt{applicability\_status} (\texttt{in\_domain}, \texttt{near\_boundary}, \texttt{out\_of\_domain}, or \texttt{calibration\_unavailable}), \texttt{nearest\_supported\_regime}, and \texttt{distance\_to\_supported\_domain}. Calibrated outputs are additionally gated by null-layer validity diagnostics (\texttt{sigma0\_valid}, \texttt{sigma0\_floored}, \texttt{fallback\_applied}). If applicability fails or null calibration is invalid/floored, then
\[
c\mathrm{EII}_{\mathrm{gene}}=\varnothing,\qquad c\mathrm{EII}_{\mathrm{site}}=\varnothing,
\]
with class labels set to \texttt{calibration\_unavailable}. Raw EII and matched-neutral significance remain available in all cases.

Calibrated decision bands are empirical and versioned (\texttt{calibration\_version}); they are not universal constants and must be re-estimated when calibration assets change.

Matched-neutral significance remains distinct from cEII:
\[
p_{\mathrm{emp}}=\frac{1+\sum_{m=1}^{M}\mathbf{1}\!\left[D_0^{(m)}\ge D_{\mathrm{obs}}\right]}{M+1},\qquad
q_{\mathrm{emp}}=\mathrm{BH}(p_{\mathrm{emp}}).
\]
The default significance decision is $\texttt{significant\_bool}=(q_{\mathrm{emp}}\le 0.05)$.

\subsection{Validation and benchmarking}

Validation used nuisance-factorial simulations with strict split separation at scenario level (\texttt{train}/\texttt{calibration}/\texttt{test}/\texttt{ood}), so threshold learning and evaluation were not performed on the same regimes. OOD scenarios were defined a priori as high-noise or high-recombination settings on deep trees.

The benchmark varied: number of taxa, alignment length, divergence/depth regime, episodic sparsity, episodic intensity, foreground-branch fraction, recombination segmentation, rate heterogeneity, codon-usage heterogeneity, and alignment noise/missingness, including neutral and near-neutral controls.

Primary evaluation included calibration/reliability metrics (Brier, ECE), ROC/PR style discrimination on $I_{\mathrm{gene}}$ and $I_{\mathrm{site}}$, threshold operating characteristics (sensitivity/specificity/PPV/NPV/FDR), and robustness/OOD stratification.

\begin{table}[t]
\centering
\caption{Core symbols and locked definitions used in release v1.1.0.}
\begin{tabular}{p{0.20\linewidth} p{0.74\linewidth}}
\hline
Symbol & Definition \\
\hline
$D_{\mathrm{obs}}$ & $\mathrm{Var}_{\mathrm{ddof}=1}(W_{0.05}(\overline{\ell}_i))$ across codon sites (locked). \\
$\mu_0,\sigma_0$ & Mean and SD of matched-neutral $D$ replicates. \\
$\mathrm{EII}_{z,\mathrm{raw}}$ & Raw standardized dispersion $(D_{\mathrm{obs}}-\mu_0)/\sigma_0$. \\
$\mathrm{EII}_{01,\mathrm{raw}}$ & $\sigma(\mathrm{EII}_{z,\mathrm{raw}})$, bounded transform of raw EII. \\
$R_{\mathrm{gene}},R_{\mathrm{site}}$ & Continuous recoverability composites from truth-aware metrics. \\
$I_{\mathrm{gene}},I_{\mathrm{site}}$ & Binary recoverability labels from pre-registered thresholds. \\
$A(x)$ & Applicability function for calibrated output validity (\texttt{in\_domain}/\texttt{near\_boundary}/\texttt{out\_of\_domain}). \\
$c\mathrm{EII}_{\mathrm{gene}},c\mathrm{EII}_{\mathrm{site}}$ & Conditional calibrated $P(I=1\mid\text{data},A=\texttt{in\_domain},\texttt{ceii\_v2})$; otherwise abstain ($\varnothing$). \\
$p_{\mathrm{emp}},q_{\mathrm{emp}}$ & Matched-neutral exceedance significance and BH-adjusted q-value. \\
\hline
\end{tabular}
\end{table}

\begin{table}[t]
\centering
\caption{Simulation benchmark factors used for empirical calibration and OOD validation.}
\begin{tabular}{p{0.34\linewidth} p{0.58\linewidth}}
\hline
Factor & Regimes/bins \\
\hline
Taxa and divergence & small/shallow, medium, large/deep trees \\
Alignment length & short, medium, long coding-region bins \\
Episodic sparsity/intensity & neutral, low, medium, high burden/effect regimes \\
Foreground-branch fraction & regime-specific random draws with fixed scenario seed \\
Recombination segmentation & none, moderate, high block-wise mixing \\
Rate heterogeneity & low, moderate, high log-normal site-rate variance \\
Alignment noise/missingness & none, mild, high nucleotide perturbation \\
Codon-usage heterogeneity & Dirichlet-derived codon preference offsets in simulator \\
Validation split & train/calibration/test + predefined OOD holdout regimes \\
\hline
\end{tabular}
\end{table}

\section{Results}

\subsection{Calibration benchmark and split structure}

The locked expanded benchmark used 80 simulated genes with strict scenario-level split separation after sigma-valid filtering: train $n=36$, calibration $n=12$, test $n=10$, and OOD $n=22$. OOD scenarios were predefined nuisance-heavy conditions (high recombination and/or high alignment noise on larger trees).

\subsection{Raw EII versus calibrated cEII}

Raw EII remained a useful continuous dispersion score but was not directly interpretable as identifiability probability. Under the revised \texttt{ceii\_v2} path, calibrated output is conditional on both applicability support and numerically valid neutral dispersion. The legacy \texttt{ceii\_v1} behavior that emitted a flat fallback value is no longer used for inference claims; in the revised runtime, unsupported inputs receive raw EII + matched-neutral significance with explicit \texttt{calibration\_unavailable} metadata rather than forced numeric cEII.

\subsection{Reliability and operating characteristics}

On held-out test data, gene-scale calibration yielded Brier $=0.291$ and ECE $=0.307$; OOD values were Brier $=0.448$ and ECE $=0.461$. Site-scale calibration was sharper in-domain (test: Brier $=0.0017$, ECE $=0.017$) but degraded out-of-domain (OOD: Brier $=0.173$, ECE $=0.208$). These quantitative results support the intended interpretation: calibration quality is regime-conditional and OOD applicability must be surfaced rather than hidden.

\subsection{Failure analysis and abstention under unsupported regimes}

Before the publication-path refactor, the expanded benchmark run using frozen-reference neutral lookup showed complete collapse of the null layer: $\texttt{fraction\_sigma0\_at\_floor}=1.0$, $\texttt{fraction\_fallback\_applied}=1.0$, and only 19/80 finite pre-floor $\sigma_0$ values. Audit of case-level failures showed two dominant modes: \texttt{missing\_neutral\_reference} (61/80) and \texttt{neutral\_group\_below\_minimum} (19/80). This confirms that the prior run operated outside valid neutral support and should not be used for calibrated cEII claims.

After replacing the force-calibration path with empirical matched-neutral calibration and abstention-aware semantics (\texttt{ceii\_v2}, winsorized dispersion, $\sigma$ floor $=10^{-3}$), the same expanded benchmark yielded $\texttt{fraction\_sigma0\_at\_floor}=0.0$ and $\texttt{fraction\_fallback\_applied}=0.0$ with full sigma-valid coverage (80/80). This behavior underpins the release interpretation: cEII is emitted only when calibration assumptions are satisfied, and abstention is an explicit safeguard against false certainty.

\subsection{Empirical applicability and abstention behavior}

On empirical panels, the release policy is now fixed: raw EII and $q_{\mathrm{emp}}$ are always reportable, but cEII is emitted only for in-domain, sigma-valid runs. Out-of-domain or numerically invalid runs return explicit \texttt{calibration\_unavailable} classes with machine-readable diagnostics. This behavior is intentional and prevents over-interpretation of unsupported calibrated probabilities.

\subsection{cEII and matched-neutral significance are complementary}

The significance layer ($p_{\mathrm{emp}}, q_{\mathrm{emp}}$) remained distinct from cEII. Inference with low cEII but significant $q_{\mathrm{emp}}$ was observed in nuisance-heavy settings, indicating excess dispersion relative to matched neutral expectation without robust localization. Conversely, moderate cEII with non-significant $q_{\mathrm{emp}}$ occurred in sparse/weak regimes, indicating partial recoverability without strong exceedance evidence. This confirms that cEII probability and q-based significance answer different questions and should be reported jointly.

\subsection{Software behavior is deterministic and robust}

Repeated runs on identical inputs produced identical outputs, indicating deterministic inference after model freezing. Invalid inputs, including empty alignments, non-codon-length sequences and single-taxon files, were rejected with explicit error messages rather than silent corruption. Conserved negative controls generated nearly flat score profiles, whereas synthetic alignments containing localized branch-specific perturbation produced corresponding peaks in inferred burden. These checks support the practical robustness of the implementation.

\section{Discussion}

BABAPPAi is built around a simple but important shift in target: from \emph{detection} of positive selection under a codon-model null to \emph{identifiability} of episodic branch--site structure from the available alignment. These goals are related but not identical. In realistic settings, biological episodic perturbation may be present while the statistical information needed to localize or test it remains weak. A method that explicitly diagnoses this limitation can therefore add value even when classical branch--site tools remain the primary inferential engines.

A second central result is that aggregation stabilizes signal. In the present benchmarks, episodic structure was more recoverable at the gene level than at individual sites. This is both biologically and statistically plausible: episodic adaptation is often sparse, lineage-restricted and transient, so local evidence may be weak even when the gene as a whole carries a real episodic burden. BABAPPAi does not eliminate this difficulty; rather, it quantifies where along the scale hierarchy recoverability begins to emerge.

The revised inference layer explicitly separates three outputs: raw EII magnitude, calibrated cEII probability, and matched-neutral significance. Raw EII is useful as a standardized dispersion effect size. cEII is the empirically calibrated probability that recoverability criteria are met in the validated simulation domain, and it is intentionally nullable outside applicability support. $q_{\mathrm{emp}}$ indicates excess branch--site dispersion relative to matched-neutral expectation. None of these outputs, alone or combined, is interpreted as direct proof of adaptive substitution.

The method is intended to complement, not replace, established branch--site procedures such as PAML and HyPhy \cite{Yang2007,Murrell2012,KosakovskyPond2011,KosakovskyPond2020}. In practice, classical likelihood-based tests can still estimate codon-model effects, while BABAPPAi quantifies whether the alignment appears informative enough for stable episodic interpretation. Agreement between strong classical evidence and high $c\mathrm{EII}_{\mathrm{gene}}$ is expected to increase confidence; disagreement flags caution and additional sensitivity analyses.

Model immutability is a design strength for reproducibility. Because BABAPPA$\Omega$ weights are fixed, methodological updates can be audited as changes in calibration, significance control, applicability envelopes, and reporting logic rather than conflated with shifting inference parameters. This separation between fixed inference and evolving interpretation is intentional scientific practice: it allows correction of overconfident post-processing without rewriting the underlying neural mapping each release.

Several limitations should be noted. First, cEII is simulator-conditional; calibration quality depends on the realism and coverage of the benchmark domain. Second, site-level identifiability remained less stable than gene-level identifiability in OOD strata, so localization claims require stronger evidence and more conservative interpretation. Third, the current benchmark used moderate sample sizes in each held-out split; wider regime coverage is still needed before claiming universal thresholds. Finally, empirical application requires careful mapping of operational branch labels onto biological hypotheses and should be accompanied by classical branch--site analyses rather than replaced by BABAPPAi outputs alone.

\section{Conclusion}

BABAPPAi reframes branch--site analysis around recoverability and matched-neutral calibration. By combining forward simulation under mutation--selection dynamics with likelihood-free neural inference, it reports raw dispersion diagnostics and empirically calibrated identifiability probabilities ($c\mathrm{EII}_{\mathrm{gene}}$, $c\mathrm{EII}_{\mathrm{site}}$) while keeping matched-neutral significance as a separate inferential layer. It does not substitute for classical branch--site testing and does not claim proof of adaptive substitution.

\section*{Availability and implementation}

BABAPPAi is available as an open-source Python package at \url{https://github.com/krishnendusinha/babappai} and as the installable package \texttt{babappai} with CLI command \texttt{babappai}. The released software reports raw EII diagnostics, conditional calibrated cEII probabilities/classes (nullable under abstention), explicit applicability diagnostics (\texttt{applicability\_score}, \texttt{applicability\_status}, nearest-regime and distance metadata), and empirical matched-neutral p/q significance outputs with calibration/model provenance fields. The canonical frozen BABAPPA$\Omega$ model artifact used by BABAPPAi is archived at \url{https://doi.org/10.5281/zenodo.18195869}; BABAPPAi software-release provenance is archived at \url{https://doi.org/10.5281/zenodo.18520163}. Benchmarking artifacts, execution logs and related reproducibility materials are archived at \url{https://doi.org/10.5281/zenodo.18197957}.

\section*{Data availability}

All benchmark data used in this study were generated by simulation. The source code, frozen inference artifacts and archived benchmark outputs required to reproduce the analyses are available through the repository and Zenodo records listed above.

\section*{Supplementary data}

Supplementary data, including additional benchmark summaries, stress-test logs and implementation details, should accompany the online version of the manuscript.

\section*{Acknowledgements}

I acknowledge the open-source scientific software ecosystem on which this work depends. I also acknowledge the personal inspiration behind the name ``BABAPPA,'' which originated from my son's imaginative word for a butterfly.

\section*{Funding}

No specific external funding was received for this work.

\section*{Conflict of interest}

The author declares no conflict of interest.

\begin{thebibliography}{99}

\bibitem{Anisimova2001}
Anisimova M, Bielawski JP, Yang Z.
Accuracy and power of the likelihood ratio test in detecting adaptive molecular evolution.
\emph{Mol Biol Evol} 2001;18:1585--1592.

\bibitem{Anisimova2003}
Anisimova M, Nielsen R, Yang Z.
Effect of recombination on the accuracy of the likelihood method for detecting positive selection at amino acid sites.
\emph{Genetics} 2003;164:1229--1236.

\bibitem{AnisimovaYang2007}
Anisimova M, Yang Z.
Multiple hypothesis testing to detect lineages under positive selection that affects only a few sites.
\emph{Mol Biol Evol} 2007;24:1219--1228.

\bibitem{Cranmer2020}
Cranmer K, Brehmer J, Louppe G.
The frontier of simulation-based inference.
\emph{Proc Natl Acad Sci USA} 2020;117:30055--30062.

\bibitem{Flagel2019}
Flagel L, Brandvain Y, Schrider DR.
The unreasonable effectiveness of convolutional neural networks in population genetic inference.
\emph{Mol Biol Evol} 2019;36:220--238.

\bibitem{FletcherYang2010}
Fletcher W, Yang Z.
The effect of insertions, deletions, and alignment errors on the branch-site test of positive selection.
\emph{Mol Biol Evol} 2010;27:2257--2267.

\bibitem{Gillespie1977}
Gillespie DT.
Exact stochastic simulation of coupled chemical reactions.
\emph{J Phys Chem} 1977;81:2340--2361.

\bibitem{HalpernBruno1998}
Halpern AL, Bruno WJ.
Evolutionary distances for protein-coding sequences: modeling site-specific residue frequencies.
\emph{Mol Biol Evol} 1998;15:910--917.

\bibitem{Hasegawa1985}
Hasegawa M, Kishino H, Yano T.
Dating the human-ape splitting by a molecular clock of mitochondrial DNA.
\emph{J Mol Evol} 1985;22:160--174.

\bibitem{JordanGoldman2012}
Jordan G, Goldman N.
The effects of alignment error and alignment filtering on the sitewise detection of positive selection.
\emph{Mol Biol Evol} 2012;29:1125--1139.

\bibitem{KosakovskyPond2011}
Kosakovsky Pond SL, Frost SDW, Muse SV.
A random effects branch-site model for detecting episodic diversifying selection.
\emph{Mol Biol Evol} 2011;28:3033--3043.

\bibitem{KosakovskyPond2020}
Kosakovsky Pond SL, Poon AFY, Velazquez R, Weaver S, Hepler NL, Murrell B, Shank SD, Magalis BR, Bouvier D, Nekrutenko A \emph{et al.}
HyPhy 2.5---a customizable platform for evolutionary hypothesis testing using phylogenies.
\emph{Mol Biol Evol} 2020;37:295--299.

\bibitem{Murrell2012}
Murrell B, Wertheim JO, Moola S, Weighill T, Scheffler K, Kosakovsky Pond SL.
Detecting individual sites subject to episodic diversifying selection.
\emph{PLoS Genet} 2012;8:e1002764.

\bibitem{SheehanSong2016}
Sheehan S, Song YS.
Deep learning for population genetic inference.
\emph{PLoS Comput Biol} 2016;12:e1004845.

\bibitem{Wong2004}
Wong WSW, Yang Z, Goldman N, Nielsen R.
Accuracy and power of statistical methods for detecting adaptive evolution in protein coding sequences and for identifying positively selected sites.
\emph{Genetics} 2004;168:1041--1051.

\bibitem{Yang2007}
Yang Z.
PAML 4: phylogenetic analysis by maximum likelihood.
\emph{Mol Biol Evol} 2007;24:1586--1591.

\bibitem{YangDosReis2011}
Yang Z, dos Reis M.
Statistical properties of the branch-site test of positive selection.
\emph{Mol Biol Evol} 2011;28:1217--1228.

\end{thebibliography}

\end{document}
